% ==============================================================================
% 00_vorspann.tex – Titelseite, Kurzfassung, Abstract, Verzeichnisse
% (Angaben zur Arbeit und das Aussehen der Titelseite stehen in preamble.tex)
% ==============================================================================

% --- Titelseite ----------------------------------------------------------------
\BHTTitelseite

% --- Vorspann: römische Seitenzahlen --------------------------------------------
\pagenumbering{Roman}

% --- Kurzfassung ---------------------------------------------------------------
\chapter*{Kurzfassung}
\addcontentsline{toc}{chapter}{Kurzfassung}
Webcam-basiertes Eye-Tracking ist kostengünstig und auf fast jedem Laptop
verfügbar. Mit allgemeinen Modellen ist es aber oft ungenau, weil sich Personen,
Kameras und Umgebungen stark unterscheiden. Diese Arbeit untersucht, wie sich
ein vortrainiertes Bildverarbeitungsmodell mit wenigen Kalibrierungsdaten an
eine einzelne Person anpassen lässt.

Als Feature Extractor dient ein auf ImageNet vortrainiertes ResNet-18, das
eingefroren bleibt. Trainiert wird nur ein kleiner Regression Head: zuerst auf
dem Datensatz GazeCapture (Basismodell) und danach mit eigenen
Webcam-Aufnahmen aus 18 Sitzungen einer Person (Personalisierung). Als
Fehlermaß dient der Diagonalfehler, also der Fehler in Prozent der
Bildschirmdiagonale.

Auf den eigenen Aufnahmen ist das Basismodell nicht besser als eine konstante
Schätzung der Bildschirmmitte (30,6\,\% gegenüber 30,3\,\%); es dient daher
vor allem als Startpunkt. Bei einer realistischen, zeitlichen Aufteilung der
Daten (Training auf dem Anfang, Test auf dem Ende einer Aufnahme) senkt die
Personalisierung den Fehler von 31,8\,\% auf 27,9\,\%, also um etwa 12\,\%.
In Sitzungen mit guten Aufnahmebedingungen ist der Gewinn deutlich größer (46
bis 47\,\%) und gilt auch für eine neue Aufnahme. Bei einer zufälligen
Aufteilung der Bilder sinkt der Fehler auf 19,8\,\% (etwa ein Drittel). Dieser
Wert ist aber zu optimistisch, weil fast gleiche Nachbarbilder in Training und
Test liegen. Ein Trainingslauf dauert auf einer CPU nur wenige Sekunden. Hintergrund und Beleuchtung haben einen größeren Einfluss als
die Gestaltung der Kalibrierung. Das Entfernen falsch gelabelter Bilder nach
Zielsprüngen verbessert das Ergebnis nur wenig.

Die Ergebnisse zeigen: Eine dateneffiziente Personalisierung ist mit einfachen
Mitteln möglich. Ihr Nutzen muss aber mit einer realistischen Aufteilung der
Daten bewertet werden. Alle Aufnahmen stammen von einer Person; die Ergebnisse
gelten daher zunächst nur für diese Person.

\Schlagwoerter{Schlagwörter}{Eye-Tracking, Blickschätzung, Webcam,
Personalisierung, Transfer Learning, ResNet-18, GazeCapture}

% --- Abstract (nur mit \ZweiteSprachetrue, siehe preamble.tex) -----------------
\ifZweiteSprache
\begin{otherlanguage}{english}
\chapter*{Abstract}
\addcontentsline{toc}{chapter}{Abstract}
Webcam-based eye tracking is inexpensive and available on almost every laptop.
With generic models, however, it is often inaccurate, because people, cameras
and environments differ strongly. This thesis investigates how a pretrained
vision model can be adapted to a single person with only a small amount of
calibration data.

A ResNet-18 pretrained on ImageNet serves as a frozen feature extractor. Only a
small regression head is trained: first on the GazeCapture dataset (base model)
and then on own webcam recordings from 18 sessions of one person
(personalization). Accuracy is measured with the diagonal error, i.e.\ the
error in percent of the screen diagonal.

On the own recordings, the base model is not better than a constant prediction
of the screen centre (30.6\,\% versus 30.3\,\%). Personalization reduces the
error to 19.8\,\%, i.e.\ by about one third. One training run takes only a few
seconds on a CPU. With a realistic, temporal split of the data the gain is
smaller (about 12\,\% compared with the constant prediction). In sessions with
good recording conditions, however, it remains clear and also transfers to a
new recording. Background and lighting have a larger influence than the design
of the calibration. Removing wrongly labelled frames after target jumps
improves the result only slightly.

The results show that data-efficient personalization is possible with simple
means. Its benefit, however, must be assessed with a realistic data split.

\Schlagwoerter{Keywords}{eye tracking, gaze estimation, webcam,
personalization, transfer learning, ResNet-18, GazeCapture}
\end{otherlanguage}
\fi

% --- Verzeichnisse -------------------------------------------------------------
\tableofcontents
\listoffigures
\listoftables

% --- Abkürzungen und Formelzeichen ---------------------------------------------
\chapter*{Abkürzungsverzeichnis}
\addcontentsline{toc}{chapter}{Abkürzungsverzeichnis}
{\renewcommand{\arraystretch}{1.08}
\begin{longtable}{@{}p{2.4cm}p{\dimexpr\textwidth-2.4cm-2\tabcolsep\relax}@{}}
  \toprule
  \textbf{Abkürzung} & \textbf{Bedeutung} \\
  \midrule
  \endhead
  \bottomrule
  \endfoot
  AdamW   & Adam-Optimierer mit entkoppeltem Weight Decay \\
  CLIP    & Contrastive Language-Image Pre-training \\
  CNN     & Convolutional Neural Network (faltendes neuronales Netz) \\
  CPU     & Central Processing Unit (Hauptprozessor) \\
  CI      & Confidence Interval (Konfidenzintervall) \\
  DA      & Domain Adaptation (Domänenadaption) \\
  DG      & Domain Generalization (Domänengeneralisierung) \\
  DINOv2  & Self-Distillation with No Labels, Version 2 \\
  FC      & Fully Connected (vollständig verbundene Schicht) \\
  fps     & Frames per Second (Bilder pro Sekunde) \\
  L1      & Loss aus dem mittleren absoluten Fehler (L1-Verlust) \\
  LoRA    & Low-Rank Adaptation \\
  MAE     & Masked Autoencoder \\
  MAML    & Model-Agnostic Meta-Learning \\
  MLP     & Multi-Layer Perceptron (mehrschichtiges Perzeptron) \\
  PEFT    & Parameter-Efficient Fine-Tuning (parametereffizientes Fine-Tuning) \\
  PoR     & Point of Regard (Blickpunkt auf dem Bildschirm) \\
  ReLU    & Rectified Linear Unit \\
  RMSE    & Root Mean Squared Error (Wurzel des mittleren quadratischen Fehlers) \\
  SD      & Standardabweichung \\
  SVR     & Support-Vector-Regression \\
  TTA     & Test-Time Adaptation (Anpassung zur Testzeit) \\
  ViT     & Vision Transformer \\
  \midrule
  \multicolumn{2}{@{}l@{}}{\textbf{Formelzeichen}} \\[0.3ex]
  $\Ediag$ & Diagonalfehler: euklidischer RMSE in Prozent der Bildschirmdiagonale \\
\end{longtable}}

% --- Verzeichnis der Fachbegriffe (eigene Datei, auch im Auszug) ---------------
% ==============================================================================
% 00b_fachbegriffe.tex – Verzeichnis der Fachbegriffe (Englisch – Deutsch)
% Wird in 00_vorspann.tex und in Abgabe.tex eingebunden.
% ==============================================================================
\chapter*{Verzeichnis der Fachbegriffe}
\addcontentsline{toc}{chapter}{Verzeichnis der Fachbegriffe}
In dieser Arbeit werden die üblichen englischen Fachbegriffe verwendet. Die
Tabelle nennt die deutsche Übersetzung und eine kurze Erklärung.

{\small\renewcommand{\arraystretch}{1.08}
\begin{longtable}{@{}L{4.1cm}L{3.8cm}>{\raggedright\arraybackslash}p{\dimexpr\textwidth-7.9cm-4\tabcolsep\relax}@{}}
  \toprule
  \textbf{Englisch} & \textbf{Deutsch} & \textbf{Kurz erklärt} \\
  \midrule
  \endhead
  \bottomrule
  \endfoot
  Accuracy & Genauigkeit & systematische Abweichung vom wahren Blickpunkt \\
  Activation (Function) & Aktivierung(sfunktion) & Funktion am Ausgang einer Schicht, z.\,B. Sigmoid \\
  Backbone & Grundnetz & großer vortrainierter Teil des Modells, der die Features berechnet \\
  Baseline & Vergleichswert & einfaches Vergleichsverfahren, z.\,B. immer derselbe Punkt \\
  Batch Size & Batchgröße & Anzahl der Bilder pro Trainingsschritt \\
  Chicken-and-Egg Problem & Henne-Ei-Problem & zwei Dinge hängen gegenseitig voneinander ab \\
  Confidence Interval (CI) & Konfidenzintervall & Bereich, in dem der wahre Mittelwert mit hoher Wahrscheinlichkeit liegt \\
  Crowdsourcing & -- & viele Freiwillige nehmen Daten selbst auf \\
  Data Leakage & Datenleck & Information aus den Testdaten gelangt ins Training \\
  Data Loss & Datenverlust & Anteil der Messungen ohne gültigen Wert \\
  Domain Adaptation & Domänenadaption & Anpassung an eine neue Umgebung mit Daten von dort \\
  Domain Generalization & Domänengeneralisierung & gute Leistung in neuen Umgebungen ohne Daten von dort \\
  Domain Shift & Domänenwechsel & Testdaten unterscheiden sich von den Trainingsdaten \\
  Early Stopping & vorzeitiges Beenden & Training endet, wenn sich die Validierung nicht mehr verbessert \\
  Epoch & Epoche & ein Durchlauf durch alle Trainingsdaten \\
  Feature(s) & Merkmal(e) & Zahlen, die ein Netz aus einem Bild berechnet \\
  Feature Extractor & Merkmalsextraktor & Netz, das die Features berechnet \\
  Feature Space & Merkmalsraum & Raum, in dem die Features liegen \\
  Fine-Tuning & Feinabstimmung & ein vortrainiertes Modell mit neuen Daten weitertrainieren \\
  Gaze Regression & Blickregression & Blick als kontinuierlicher Wert vorhersagen \\
  Head-Training & -- & nur der Regression Head wird trainiert \\
  Image / Text Encoder & Bild- / Textencoder & Netz, das ein Bild bzw. einen Text in Features umwandelt \\
  Label Noise & Label-Rauschen & Labels stimmen nicht genau mit dem echten Blick überein \\
  Learning Rate & Lernrate & Schrittweite beim Training \\
  Loss / Loss Function & Verlust / Verlustfunktion & Maß für den Fehler, das beim Training minimiert wird \\
  Overfitting & Überanpassung & Modell lernt die Trainingsdaten auswendig und wird auf neuen Daten schlechter \\
  Precision & Präzision & Streuung wiederholter Messungen \\
  Pretraining & Vortraining & Training auf sehr vielen allgemeinen Daten vor der eigentlichen Aufgabe \\
  Regression Head & Regressionskopf & kleines Netz, das aus den Features die Koordinaten berechnet \\
  Remote Camera & entfernte Kamera & Kamera steht vor der Person, nicht am Kopf \\
  Ridge Regression & Ridge-Regression & lineare Regression mit Strafe für große Gewichte \\
  Saccade & Sakkade & sehr schnelle Augenbewegung zu einem neuen Ziel \\
  Self-Supervised Learning & selbstüberwachtes Lernen & Lernziel entsteht aus den Daten selbst, ohne Labels \\
  Split & Aufteilung & Aufteilung der Daten in Training, Validierung und Test \\
  Supervised Learning & überwachtes Lernen & Lernen mit Labels, die Menschen vergeben haben \\
  Test-Time Adaptation (TTA) & Anpassung zur Testzeit & Anpassung erst während der Nutzung \\
  Transition Frames / Period & Übergangsbilder / Übergangszeit & Bilder bzw. Zeit direkt nach einem Zielsprung \\
  Unconstrained Conditions & uneingeschränkte Bedingungen & freie Kopfhaltung, wechselndes Licht, beliebiger Abstand \\
  Unsupervised Domain Adaptation & unüberwachte Domänenadaption & Domain Adaptation ohne Labels der Zielumgebung \\
  Weakly Supervised & schwach überwacht & Lernen mit wenigen oder groben Labels \\
  Weight Decay & Gewichtsabnahme & Strafe für große Gewichte gegen Overfitting \\
\end{longtable}}


% --- Hauptteil: arabische Seitenzahlen ------------------------------------------
\cleardoublepage
\pagenumbering{arabic}
